% Cài đặt Odoo 19 (dev) — compile: lualatex/xelatex odoo19-setup.tex
\documentclass[11pt,a4paper]{article}
\usepackage{fontspec}
\setmainfont{Lato}
\setmonofont{DejaVu Sans Mono}[Scale=0.9]
\usepackage[margin=2.3cm]{geometry}
\usepackage{parskip}
\usepackage{enumitem}
\usepackage{listings}
\usepackage{xcolor}
\usepackage[hidelinks]{hyperref}
\definecolor{codebg}{RGB}{245,245,248}
\lstset{basicstyle=\ttfamily\small, backgroundcolor=\color{codebg},
  breaklines=true, columns=fullflexible, frame=single, framesep=4pt,
  rulecolor=\color{black!15}, xleftmargin=4pt}

\begin{document}

\begin{center}{\LARGE\bfseries Cài đặt Odoo 19 (bản dev / thuần)}\end{center}
\vspace{4pt}

\section*{1. Cần cài sẵn}
\begin{itemize}[nosep]
  \item \textbf{git}, \textbf{Miniconda} (Python env), \textbf{PostgreSQL} 14+.
  \item \textbf{wkhtmltopdf} bản Qt-patched 0.12.6 (chỉ để in PDF; thiếu vẫn chạy được) --- tải ở \url{https://github.com/wkhtmltopdf/packaging/releases}.
  \item Node đi kèm trong \texttt{environment.yml} (để cài \texttt{rtlcss}).
\end{itemize}

\section*{2. Nguồn Odoo 19}
Repo chính thức: \url{https://github.com/odoo/odoo} --- nhánh \texttt{19.0}.

\section*{3. Thư viện Python}
Dùng file \texttt{environment.yml} đi kèm (tạo env conda \texttt{odoo19} + cài toàn bộ package Odoo 19).

\section*{4. Các lệnh chạy}
Chạy lần lượt 5 nhóm: (1) thư viện hệ thống, (2) clone Odoo, (3) môi trường Python,
(4) PostgreSQL + tạo role, (5) chạy Odoo. Xong mở \url{http://localhost:8069}
(lần đầu tạo DB qua web).

\begin{lstlisting}
sudo apt update && sudo apt install -y git build-essential python3-dev \
  libpq-dev libxml2-dev libxslt1-dev libldap2-dev libsasl2-dev \
  libjpeg-dev zlib1g-dev libssl-dev libffi-dev

git clone --depth 1 --branch 19.0 https://github.com/odoo/odoo.git ~/odoo

conda env create -f environment.yml
conda activate odoo19
npm install -g rtlcss

sudo apt install -y postgresql
sudo service postgresql start
sudo -u postgres createuser -s $USER

cd ~/odoo
python odoo-bin --addons-path=addons -d odoo_dev
\end{lstlisting}

\section*{5. (Tuỳ chọn) File config}
Mục 4 chạy không cần config. Khi bắt đầu viết addon riêng + muốn hot-reload thì tạo
file \texttt{odoo.conf} (đổi \texttt{USER} thành user máy bạn):

\begin{lstlisting}
[options]
addons_path = /home/USER/odoo/addons,/home/USER/odoo-custom
data_dir = /home/USER/odoo/.local-data
db_host = 127.0.0.1
db_port = 5432
db_user = USER
http_port = 8069
workers = 0
admin_passwd = admin
\end{lstlisting}

Chạy với config (kèm hot-reload):
\begin{lstlisting}
python odoo-bin -c odoo.conf -d odoo_dev --dev=all
\end{lstlisting}

\end{document}
